\subsection{Komplexe Differenzierbarkeit}
\begin{emphbox}
    Sei $U\subset\mathbb C$ offen. Ein $f:U\rightarrow\mathbb C$ ist komplex differenzierbar (\textbf{holomorph}) bei $z_0\in U$ g.d.w. folgender Grenzwert existiert
    \begin{equation*}
        \lim_{z\rightarrow z_0}\frac{f(z) - f(z_0)}{z-z_0}=:f'(z_0)
    \end{equation*}

    Sei $U$ ein einfach zusammenhängendes Gebiet und $z_0\in U$ bel. Sei $f:U\rightarrow\mathbb C$ holomorph auf $U$, dann ist
    \begin{equation*}
        F(z) = \int_\gamma f(\xi)\diff\xi
    \end{equation*}
    die Stammfunktion von $f$, wobe $\gamma[a,b]\rightarrow U$ eine bel. Kurve mit $\gamma(a) = z_0, \gamma(b) = z$
\end{emphbox}
Falls $U=\mathbb C$ ist $f$ eine ganze Funktion

\begin{sectionbox}
    \subsubsection{Cauchy-Riemann Differentialgleichungen}
    Sei $f(x+iy) = u(x,y) + iv(x,y)$ mit $u,v:\mathbb R^2 \rightarrow \mathbb R$, dann ist äquivalent
    \begin{itemize}
        \item $f$ ist bei $z_0 = x_0+iy_0$ komplex differenzierbar
        \item Die Funktion $(x,y)\mapsto \begin{pmatrix}u(x,y) \\ v(x,y)\end{pmatrix}$ ist bei $(x_0,y_0)\in\mathbb R^2$ (reell) differenzierbar und $\partial_x u = \partial_y v, \qquad \partial_y u = -\partial_x v$
        \item Es gilt die Ableitung $f'(x+iy) = \partial_x u(x,y) - i\partial_y u(x,y)=\partial_y v(x,y) + i\partial_x v(x,y)$
    \end{itemize}
\end{sectionbox}
\begin{sectionbox}
    \subsubsection{Wichtige Sätze}
    \begin{itemize}
        \item Sei $U$ offen und zusammenhängend und $f$ holomorph auf $U$ und $C\ge 0$ so, dass $|f(z)|=C \; \forall z\in U \Rightarrow f$ konstant auf $U$
    \end{itemize}
\end{sectionbox}

\subsection{Potenzreihen in $\mathbb C$}
\begin{emphbox}
    \begin{equation*}
        P(z) = \sum_{n=1}^{\infty}a_n(z-z_0)^n
    \end{equation*}
\end{emphbox}
\begin{sectionbox}
    \subsubsection{Konvergenz}
    \begin{center}
        \boxed{R = \lim\limits_{n\rightarrow \infty} \abs{\frac{a_n}{a_{n+1}}}=\frac{1}{\lim\limits_{n\rightarrow \infty}\sqrt[n]{\abs{a_n}}}} \\
        \boxed{R = \liminf\limits_{n\rightarrow \infty} \abs{\frac{a_n}{a_{n+1}}}=\frac{1}{\limsup\limits_{n\rightarrow \infty}\sqrt[n]{\abs{a_n}}}}
    \end{center}

    \begin{itemize}
        \item $P(z)$ konvergiert absolut für $z\in B_R(z_0) = \{z\in\mathbb C \, | \, |z-z_0| < R \}$
        \item $P(z)$ divergiert für $z\notin \overline{B_R(z_0)} = \{z\in\mathbb C \, | \, |z-z_0| \le R \}$
        \item $\forall \, 0 \le r < R$ ist die Konvergenz gleichmäßig. $P(z)$ ist eine stetige Funktion auf $B_R(z_0)$
    \end{itemize}
\end{sectionbox}

\begin{sectionbox}
    \subsubsection{Wichtige Sätze}
    \begin{itemize}
        \item $f$ heißt analytisch in $z_0\in U$ g.d.w. eine Potenzreihe $P(z)$ mit Konvergenzradius $R>0$ existiert, sodass $f(z)=P(z) \, \forall z\in B_R(z_0) \cap U$ \\
              Jede analytische Funktion lässt sich also lokal in eine Potenzreihe entwickeln
    \end{itemize}

    \subsubsection*{Hauptzweig des Logarithmus}
    $\log : \mathbb C\setminus\{0\} \rightarrow \mathbb R + i(-\pi, \pi], \, \log(z) = \log|z| + i\arg{z}$
\end{sectionbox}

\begin{sectionbox}
    \subsubsection{Wichtige Reihen}
    \begin{tabular}{ll}
        \(\displaystyle \frac{1}{\omega - z} =  \sum_{k=0}^\infty \frac{(z-z_0)^k}{(\omega-z_0)^{k+1}}\) & \(\displaystyle e^z=\sum_{k=0}^\infty \frac{z^k}{k!} \)                  \\
        \(\displaystyle \sin(z) = \sum_{k=0}^\infty \frac{(-1)^k}{(2k+1)!}z^{2k+1} \)                    & \(\displaystyle \cos(z) = \sum_{k=0}^\infty\frac{(-1)^k}{(2k)!}z^{2k} \)
    \end{tabular}
\end{sectionbox}


\subsection{Komplexe Kurvenintegrale}
\paragraph{geschlossen} $\gamma:[a,b]\rightarrow\mathbb C, \gamma(a)=\gamma(b)$
\paragraph{einfach geschlossen} $\gamma(s)=\gamma(t) \Leftrightarrow s=t$ oder $\{s,t\} = \{a,b\}$
\begin{sectionbox}
    \begin{emphbox}
        Sei $\gamma [a,b] \rightarrow\mathbb C$ stetig differenzierbar und $f:\text{im}(\gamma) \rightarrow\mathbb C$ stetig
        \begin{equation*}
            \int_\gamma f(z)\diff z = \int_a^bf(\gamma(t))\gamma'(t)\diff t
        \end{equation*}

        Falls $\gamma$ stückweise stetig differenzierbar mit $m$ Sprungstellen
        \begin{equation*}
            \int_\gamma f(z)\diff z = \sum_{j=0}^{m}\int_{\gamma_j} f(z)\diff z
        \end{equation*}
    \end{emphbox}

    \subsubsection{Eigenschaften}
    \begin{itemize}
        \item Ist $\hat \gamma(t) = \gamma(b-(t-a))$ die Kurve $\gamma$ umgekehrt durchlaufen, dann gilt: $\int_{\hat \gamma}f(z)\diff z = -\int_\gamma f(z)\diff z$
        \item $\oint_{|z-z_0|=r} f(z)\diff z=\oint_\gamma f(z)\diff z$ mit $\gamma(t) = z_0 + re^{it}$ ist die Kurve mit Mittelpunkt $z_0$, die einen Kreis mit Radius $r$ durchläuft
        \item $\int_{[z_0,z_1]}f(z)\diff z = \int_\gamma f(z)\diff z$ mit $\gamma(t) = z_0 + t(z_1-z_0)$ ist die Kurve, welche die Gerade zwischen $z_0$ und $z_1$ beschreibt
    \end{itemize}

    \paragraph{Bogenlänge} $L(\gamma) = \int_a^b |\gamma'(t)|\diff t$

    \paragraph{Aufteilung in Real- und Imaginärteil} Sei $\gamma:[a,b]\rightarrow\mathbb C$ stückweise $\mathcal C^1$.
    Sei $f:\text{im}(\gamma)\rightarrow\mathbb C$. Die Kurve $\tilde\gamma:[a,b]\rightarrow\mathbb R^2, \tilde\gamma(t) = \vect{\Re{\gamma(t)} \\ \Im{\gamma(t)}}$ und $f(x+iy) = u(x,y) + iv(x,y)$.
    Dann gilt \begin{align*}
        \Re{\int_\gamma f(z)\diff z} = \int_{\tilde\gamma}\vect{u \\ -v}\cdot \diff\vec x \\
        \Im{\int_\gamma f(z)\diff z} = \int_{\tilde\gamma}\vect{v \\ u}\cdot \diff\vec x
    \end{align*}
    Für ein Vektorfeld $w(x,y) = \vect{u(x,y) \\ -v(x,y)}$ ist der Realteil des Kurvenintegrals gleich der Zirkulation von $w$ entlang $\tilde\gamma$ und der Imaginärteil gleich dem Fluss von $w$ senkrecht zu $\tilde\gamma$
\end{sectionbox}

\subsection{Integralsatz von Cauchy}
\paragraph{$U\subset\mathbb C$ wegzusammenhängend} Falls $\forall z_0,z_1\in U$ die Verbindungslinie $[z_0,z_1]$ existiert
\paragraph{einfach wegzusammenhängend} Falls jede geschlosse Kurve in $U$ sich zu einem Punkt zusammenziehen lässt

\begin{emphbox}
    Sei $U$ ein einfach zusammenhängendes Gebiet und $f:U\rightarrow\mathbb C$ in $U$ komplex diffb.
    Sei $\gamma:[a,b]\rightarrow\mathbb C$ eine in $U$ verlaufende einfach geschlossene Kurve.
    Dann gilt \begin{equation*}
        \oint_\gamma f(z)\diff z = 0
    \end{equation*}
\end{emphbox}
